% Build: cd paper && latexmk -pdf main.tex
% Every table under tables/ and every number quoted in the text is generated by
% scripts/make_tables.py from results/holdout_2026. Do not edit them by hand.
\documentclass[11pt]{article}

\usepackage[margin=1in]{geometry}
\usepackage{booktabs}
\usepackage{graphicx}
\usepackage{amsmath}
\usepackage{amssymb}
\usepackage[hidelinks]{hyperref}
\usepackage{natbib}
\usepackage{microtype}

% Every quoted number is a macro generated from the results, so no sentence can
% drift from the table above it between one re-run and the next.
% Generated by scripts/make_tables.py. Do not edit.
\newcommand{\MaxGain}{28\%}
\newcommand{\MaxGainGroup}{wikipedia (weekly)}
\newcommand{\NumGroups}{7}
\newcommand{\NumGroupsFoundationBest}{6}
\newcommand{\NumGroupsFoundationBestSig}{5}
\newcommand{\NumGroupsNoDifference}{1}
\newcommand{\NumArimaGroups}{7}
\newcommand{\NumArimaBeaten}{7}
\newcommand{\ArimaMedianSlowdown}{54}
\newcommand{\ArimaMaxSlowdown}{2,689}
\newcommand{\ArimaMaxGroup}{weather (hourly)}
\newcommand{\ArimaMinSlowdown}{5}

% Generated by scripts/corpus_familiarity.py. Do not edit.
\newcommand{\CorpusDeltaWiki}{-0.53}
\newcommand{\CorpusDeltaOther}{-0.09}
\newcommand{\CorpusP}{$<10^{-5}$}
\newcommand{\CorpusEffect}{-0.13}
\newcommand{\CorpusNWiki}{1,500}
\newcommand{\CorpusNOther}{755}


\title{A Later Test Set Is Not a New Domain:\\
       Pretraining Familiarity Survives a Contamination-Free Hold-Out}
\author{Mahdi Naser Moghadasi}
\date{\today}

\begin{document}
\maketitle

\begin{abstract}
Time-series foundation models are evaluated almost exclusively on public
archives that predate them, so a strong score cannot be separated from having
seen the test set during pretraining. The obvious remedy is a hold-out that
postdates the models. We build one: thirteen forecasters --- four classical,
three trained per dataset, six pretrained --- on seven groups drawn from five
domains, every observation published after the last model was released, and
every dataset rebuildable without an API key. Under this protocol pretrained
models win \NumGroupsFoundationBestSig{} of \NumGroups{} groups, lose one to a
Theta baseline, and on daily exchange rates are indistinguishable from a
seasonal naive forecast, along with every other method tested.

We then ask what separates the wins from the losses, and report a negative
result: the two intrinsic properties one would reach for --- seasonal strength
and spectral entropy, measured on the input window --- do not account for the
pattern, and seasonal strength is if anything negatively associated with the
advantage. What does track it is corpus familiarity. Our largest gain
(\MaxGain{} lower MASE than the best classical method, on \MaxGainGroup{})
falls on Wikipedia pageviews, the domain TimesFM's authors describe as the bulk of its pretraining
corpus, at the same granularities and differing only in time window. Within the
pretrained family, where every model forecasts identical series so that series
difficulty cancels, the TimesFM family outranks the Chronos family by
\CorpusDeltaWiki{} on Wikipedia against \CorpusDeltaOther{} everywhere else
(\CorpusNWiki{} vs.\ \CorpusNOther{} series, Mann--Whitney $p$~\CorpusP{}). We
conclude that a temporal hold-out removes memorisation of a window but not
familiarity with a domain, that benchmarks therefore need domain hold-outs
stated relative to disclosed corpora, and that the practitioner's question is
less which model is better than whether their domain is one the model was
raised on.
\end{abstract}

\section{Introduction}

The claim attached to time-series foundation models is that a single
pretrained network forecasts unseen series without being fitted to them, and
that it does so better than methods estimated on the target series itself.
That claim is usually supported by evaluation on a standard collection of
public archives. Those archives are the difficulty. Every widely used
benchmark set --- the M-competition data, ETT, the Monash repository ---
predates the models being evaluated and is distributed in the same public form
that pretraining corpora are assembled from. A high score is therefore
consistent with two very different explanations, generalisation and
memorisation, and the published evaluations cannot distinguish them.

The remedy is not a better statistic but a better hold-out. This paper
evaluates on observations that did not exist when the models were trained. We
assemble seven forecasting groups from five domains, every one of them
published continuously by an institutional source, and place the entire test
window after the release date of the most recently released model in the
comparison. Nothing in the hold-out can have been memorised because none of it
had happened yet.

Under that protocol the headline finding of the foundation-model literature
survives only in part. Pretrained models are clearly better on some domains and
clearly not better on others. The interesting question is what separates them,
and our answer is not the one we set out to give: we expected a property of the
series and found a property of the corpus. The domain on which pretrained
models gain most is the domain one of them was largely pretrained on --- the
same source and the same granularities, differing only in time window --- and
the advantage is largest for exactly the model whose corpus it is.

This is not memorisation of the test set; the hold-out makes that impossible.
It is the weaker but more durable form of the same problem. A model that has
read a decade of Wikipedia pageview series has learned what Wikipedia pageview
series do, and January 2026 pageviews still do it. Moving the test window
forward controls for the first kind of contamination and leaves the second
untouched, which means the field's evaluations cannot be repaired by
recency alone.

\paragraph{Contributions.}
\begin{enumerate}
  \item \textbf{A contamination-free protocol}, with a hold-out that postdates
        every evaluated model, over five domains built only from keyless,
        continuously published sources so that any reader can rebuild the exact
        panel from the fetchers we release.
  \item \textbf{Two negative results.} Pretrained models do not beat classical
        baselines everywhere: we identify one domain where a Theta baseline
        ranks first and one where nothing beats seasonal naive. And the
        intrinsic explanation fails --- neither seasonal strength nor spectral
        entropy accounts for where the advantage appears.
  \item \textbf{Evidence that pretraining familiarity persists past a temporal
        hold-out}, from a within-family comparison on identical series that no
        property of those series can explain.
  \item \textbf{Significance rather than decimal places.} We report
        multiple-comparison rank intervals and corrected paired tests, and show
        that most published-style gaps between leading pretrained models fall
        inside them.
  \item \textbf{Cost alongside accuracy}, measured on identical hardware in the
        same runs.
\end{enumerate}

\section{Related work}

\paragraph{Pretrained forecasters.}
Chronos \citep{chronos} tokenises scaled series and trains a language-model
architecture over the token sequence; Chronos-Bolt and Chronos-2 are its
faster successors. TimesFM \citep{timesfm} is a decoder-only patched
transformer trained on a large corpus dominated by web and search traffic.
Moirai \citep{moirai} targets arbitrary frequency and multivariate input via
masked patch prediction. These models are trained on differing and only
partially disclosed corpora, which is exactly why an evaluation cannot verify
by inspection that a given test set was excluded.

\paragraph{Benchmarks and their contamination.}
The M4 \citep{m4} and M5 \citep{m5} competitions established the evaluation
conventions this paper follows: MASE and sMAPE against a seasonal naive scale,
and multiple-comparison rank intervals rather than raw mean losses. Their data,
along with the Monash archive \citep{monash}, is public and static, which makes
it ideal for comparability and unusable for testing memorisation. Recent work
has begun to raise leakage explicitly; our contribution is to make the hold-out
temporal and the collection reproducible rather than to audit corpora
post hoc.

\paragraph{Classical baselines.}
Theta \citep{theta} won M3 and remains difficult to beat on seasonal series;
automatic ARIMA and ETS \citep{hyndman2008} are the standard automated
alternatives. We use the \texttt{statsforecast} implementations so that the
baselines are the ones practitioners would actually run.

\section{Protocol}

\subsection{Contamination-free hold-out}

Let $t^\ast$ be one month after the latest disclosed pretraining cutoff among
the evaluated models. Every test observation is drawn from after $t^\ast$; the
forecast origin is fixed at 1 January 2026 for all groups, and the input
context is whatever precedes it. The cutoffs we could establish from model
cards and papers are tabulated in the supplement, together with the cases where
a cutoff is not disclosed at all --- for those we take the model's public
release date, which is conservative in the direction that favours the model.

\subsection{Domains}

All five domains are published continuously by an institution, are retrievable
without registration or an API key, and are not redistributions of an existing
benchmark. Table~\ref{tab:mase} lists the resulting groups.

\begin{description}
  \item[Wikipedia pageviews] (daily, weekly, monthly). Human attention:
        strongly periodic, heavy-tailed, punctuated by exogenous spikes.
  \item[Weather] (hourly, ERA5 reanalysis via Open-Meteo). Smooth, dominated
        by a clean daily cycle.
  \item[Air quality] (hourly, CAMS via Open-Meteo). Same sampling rate and
        daily cycle as weather, but spiky and heavy-tailed. Included
        specifically to separate ``good at hourly data'' from ``good at smooth
        data''.
  \item[Electricity] (hourly, Danish grid settlement). The domain most often
        claimed in the foundation-model literature, but evaluated here without
        the contamination that the ETT datasets carry.
  \item[Exchange rates] (daily, ECB reference rates). The adversarial case:
        the standing result in forecasting is that nothing reliably beats a
        naive forecast, so this domain tests whether a benchmark can detect the
        absence of an effect.
\end{description}

Series selection is specified before any result is seen --- the full city list,
the full currency list, the publisher's own column set --- and nothing is added
or removed afterwards.

\subsection{Metrics and significance}

We report MASE and sMAPE as defined by the M4 organisers, and weighted quantile
loss and 80\% interval coverage for the probabilistic models. Because a table
of mean losses invites ranking by the third decimal place, our primary
comparison is by rank: series are ranked across models, and each model receives
an average rank with a Nemenyi interval derived from the Friedman statistic,
following the multiple-comparisons-with-the-best presentation used to report
M5. We additionally test each model against the per-group winner with a
Wilcoxon signed-rank test, Holm-corrected across comparisons. We use
signed-rank rather than a paired $t$-test because MASE across a panel is
strongly right-skewed, and we avoid the label ``Diebold--Mariano'' because that
test compares errors over time within one series, which is not the comparison a
panel benchmark makes.

\section{Results}

\subsection{Accuracy}

\begin{table}[t]
\centering
\small
\begin{tabular}{lrrrrrrr}
\toprule
Model & Wikipedia (daily) & Wikipedia (weekly) & Wikipedia (monthly) & Weather (hourly) & Air quality (hourly) & Electricity (hourly) & Exchange rates (daily) \\
\midrule
\multicolumn{8}{l}{\emph{Classical}} \\
\quad SeasonalNaive & 1.080 & 1.642 & 0.822 & 1.359 & 1.142 & 1.412 & 2.405 \\
\quad Theta & 0.961 & 1.499 & 0.746 & 1.280 & 1.120 & \textbf{1.023} & 2.448 \\
\quad AutoETS & 1.149 & 1.817 & 0.941 & 1.558 & 1.514 & 1.239 & 2.468 \\
\quad AutoARIMA & 0.928 & 1.505 & 0.853 & 1.191 & 1.024 & 1.158 & 2.502 \\
\multicolumn{8}{l}{\emph{Trained per dataset}} \\
\quad LightGBM & 1.162 & 3.772 & 1.545 & 1.118 & 0.983 & 1.285 & 3.227 \\
\quad LSTM & 0.915 & 1.202 & 0.786 & 1.269 & 0.988 & 9538775.181 & 2.859 \\
\quad NBEATS & 0.870 & 1.326 & 0.819 & 1.093 & 0.903 & 1308279.661 & 2.539 \\
\multicolumn{8}{l}{\emph{Pretrained}} \\
\quad ChronosBoltSmall & 0.790 & 1.132 & 0.641 & 1.067 & 0.923 & 1.176 & 2.614 \\
\quad ChronosBoltBase & 0.774 & 1.143 & 0.635 & 1.055 & 0.915 & 1.254 & 2.501 \\
\quad Chronos2 & 0.734 & 1.106 & \textbf{0.620} & \textbf{1.043} & \textbf{0.895} & 1.131 & \textbf{2.390} \\
\quad TimesFM & 0.695 & 1.089 & 0.621 & 1.073 & 0.903 & 1.186 & 2.400 \\
\quad TimesFM3 & \textbf{0.689} & \textbf{1.073} & 0.623 & 1.069 & 0.901 & 1.145 & 2.418 \\
\quad Moirai2 & 0.748 & 1.135 & 0.629 & 1.067 & 0.906 & 5.326 & 2.517 \\
\bottomrule
\end{tabular}

\caption{Mean MASE by group; lower is better, best per column in bold.}
\label{tab:mase}
\end{table}

\begin{table}[t]
\centering
\small
\begin{tabular}{lrrrrrrr}
\toprule
Model & Wikipedia (daily) & Wikipedia (weekly) & Wikipedia (monthly) & Weather (hourly) & Air quality (hourly) & Electricity (hourly) & Exchange rates (daily) \\
\midrule
Chronos2 & 5.29 & \underline{5.31} & \underline{5.20} & \underline{5.48} & \underline{5.36} & \underline{5.15} & \underline{5.76} \\
TimesFM3 & \underline{4.58} & \underline{4.89} & \underline{5.44} & \underline{5.88} & \underline{5.29} & \underline{5.30} & \underline{6.55} \\
TimesFM & \underline{4.64} & \underline{4.98} & \underline{5.22} & \underline{6.17} & \underline{5.48} & \underline{6.21} & \underline{5.62} \\
Moirai2 & \underline{5.15} & 5.83 & \underline{5.15} & \underline{5.96} & \underline{5.36} & \underline{5.91} & \underline{6.24} \\
ChronosBoltBase & 6.10 & 6.00 & \underline{5.53} & \underline{5.89} & \underline{5.89} & 6.45 & \underline{6.34} \\
ChronosBoltSmall & 6.31 & 6.21 & 6.16 & \underline{6.10} & 6.03 & \underline{6.09} & \underline{7.59} \\
NBEATS & 7.24 & 7.15 & 7.59 & 6.31 & \underline{5.73} & 7.26 & \underline{6.24} \\
Theta & 7.90 & 7.51 & 7.09 & 8.18 & 7.59 & \underline{4.53} & \underline{6.00} \\
LSTM & 8.31 & 6.69 & 7.73 & 8.24 & 6.92 & 8.06 & 8.14 \\
AutoARIMA & 8.03 & 8.21 & 8.50 & 7.45 & -- & -- & -- \\
AutoETS & 9.31 & 8.79 & 8.01 & 9.46 & 9.38 & 6.70 & \underline{5.91} \\
SeasonalNaive & 9.13 & 8.24 & 8.30 & 9.18 & 8.28 & 9.43 & \underline{5.98} \\
LightGBM & 9.01 & 11.19 & 11.08 & 6.71 & 6.70 & 6.91 & \underline{7.62} \\
\bottomrule
\end{tabular}

\caption{Average rank across series (lower is better). Underlined entries are
statistically indistinguishable from the best model in that column under the
Nemenyi interval.}
\label{tab:ranks}
\end{table}

Table~\ref{tab:ranks} is the paper's central result, and it does not read like
the tables in the papers introducing these models. A pretrained model holds the
top rank in \NumGroupsFoundationBest{} of the \NumGroups{} groups, but that
count overstates the case and we do not use it: in
\NumGroupsNoDifference{} of those the Friedman test does not reject, so the
column has a winner in the arithmetic sense only. Counting a rank-first under a
null result as a win is precisely the overclaim this paper exists to argue
against. The defensible statement is that pretrained models win
\NumGroupsFoundationBestSig{} groups, lose one, and are indistinguishable from
everything else --- seasonal naive included --- on the last.

The two exceptions are the informative ones. On electricity Theta ranks first
and the leading pretrained models are merely indistinguishable from it. On
exchange rates the tied set contains every method tested; the benchmark detects
no difference between any of the thirteen, which is the right answer for that
domain rather than a failure of the design.

\subsection{The intrinsic explanation, and why we abandoned it}

Ordering the groups by the gap between the best pretrained model and the best
classical one gives a gradient: largest on Wikipedia traffic, smaller on
weather and air quality, negative on electricity, null on exchange rates. The
natural hypothesis is that some property of the series explains it, and the
natural candidates are periodicity and noise --- pretraining should help most
where a repeating structure exists but is hard to estimate from a few hundred
observations alone.

We tested that directly. For each group we measured, on the input window only,
the median seasonal strength (from an STL decomposition at the group's seasonal
period) and the median spectral entropy of the differenced series, both as
defined in the feature literature used to characterise the M4 panel
(Table~\ref{tab:features}). Neither explains the pattern. Seasonal strength is
if anything \emph{negatively} associated with the advantage: Wikipedia series
have almost no measurable seasonality at the tested period and show the largest
gains, while weather is the most strongly seasonal group in the study and shows
a middling one. Spectral entropy fails in the other direction: exchange rates
are as high-entropy as Wikipedia traffic and show no advantage at all. With
seven groups no correlation here is statistically meaningful in any case, and
we report the measurement rather than a story fitted to it.

\begin{table}[t]
\centering
\small
\begin{tabular}{lrrr}
\toprule
Group & Seasonal strength & Spectral entropy & Pretrained gain \\
\midrule
Wikipedia (daily) & 0.10 & 0.93 & +25.8\% \\
Wikipedia (weekly) & 0.00 & 0.92 & +28.4\% \\
Wikipedia (monthly) & 0.06 & 0.87 & +16.9\% \\
Weather (hourly) & 0.77 & 0.69 & +12.5\% \\
Air quality (hourly) & 0.34 & 0.84 & +12.6\% \\
Electricity (hourly) & 0.47 & 0.83 & -10.6\% \\
Exchange rates (daily) & 0.00 & 0.93 & n.s. \\
\bottomrule
\end{tabular}

\caption{Series properties measured on the input window only, against the
observed advantage. Neither column orders the last one.}
\label{tab:features}
\end{table}

\subsection{Corpus familiarity survives the hold-out}

What does line up with the gradient is not a property of the series but of the
models' training data. TimesFM's authors describe its pretraining corpus as
dominated by Wikipedia pageviews --- on the order of $10^{11}$ time points at
hourly, daily, weekly and monthly granularity, running to late 2023. Our three
Wikipedia groups are pageviews at daily, weekly and monthly granularity,
beginning in 2026. The hold-out is clean in time and identical in kind.

A domain being both the largest corpus component and the largest measured gain
is suggestive, but it is one coincidence and admits a benign reading: perhaps
Wikipedia traffic is simply a domain where all pretrained models do well. We
can separate the two, because the pretrained models forecast the \emph{same
series}. Any property of those series --- difficulty, seasonality, noise ---
applies equally to all of them and cancels in a within-family comparison. What
does not cancel is what each model read during pretraining.

Ranking the six pretrained models against each other on every series, the
TimesFM family sits \CorpusDeltaWiki{} ranks below the Chronos family on
Wikipedia and only \CorpusDeltaOther{} below it on the other four domains
(\CorpusNWiki{} vs.\ \CorpusNOther{} series; Mann--Whitney $p$~\CorpusP{},
rank-biserial \CorpusEffect{}). The model with Wikipedia in its corpus is
specifically better at Wikipedia, and only there. The effect is modest in size,
and we present it as such --- but its direction is not something the series can
account for.

We therefore read the gradient as a corpus-overlap effect rather than a
capability one, and note that this is a hypothesis our design supports rather
than a demonstrated mechanism. The decisive experiment is a benchmark whose
domains are stratified by membership in each model's disclosed corpus, which
requires disclosure the field does not currently provide. That, rather than a
further increase in the size of the pretraining set, seems to us the useful
next contribution.

\subsection{What this means for evaluation}

A temporal hold-out is necessary and not sufficient. It rules out the model
having seen these observations; it does not rule out the model having spent its
training on this kind of series, and the second effect is large enough to
reorder a leaderboard. Reporting a single average across domains hides it
entirely, because the average is dominated by whichever domains happen to be in
the benchmark, and those are the well-published ones --- which are also the
ones most likely to be in a pretraining corpus.

\subsection{Cost}

\begin{table}[t]
\centering
\small
\begin{tabular}{lrr}
\toprule
Model & Median s/series & Relative to ChronosBoltSmall \\
\midrule
SeasonalNaive & 0.0000 & free \\
ChronosBoltSmall & 0.0125 & 1.0$\times$ \\
Moirai2 & 0.0141 & 1.1$\times$ \\
ChronosBoltBase & 0.0174 & 1.4$\times$ \\
Chronos2 & 0.0181 & 1.4$\times$ \\
TimesFM3 & 0.0269 & 2.2$\times$ \\
Theta & 0.0429 & 3.4$\times$ \\
LightGBM & 0.0550 & 4.4$\times$ \\
NBEATS & 0.0627 & 5.0$\times$ \\
TimesFM & 0.0726 & 5.8$\times$ \\
LSTM & 0.0887 & 7.1$\times$ \\
AutoETS & 0.1362 & 11$\times$ \\
AutoARIMA & 3.8484 & 308$\times$ \\
\bottomrule
\end{tabular}

\caption{Median wall-clock seconds per series on identical hardware, including
model load.}
\label{tab:cost}
\end{table}

The accuracy differences among the leading pretrained models are small enough
to sit inside their rank intervals in most groups. Their costs are not
comparable in the same way. An automatic ARIMA search costs a median
\ArimaMedianSlowdown{}$\times$ more per series than a pretrained forward pass,
and that factor is not a constant: it runs from \ArimaMinSlowdown{}$\times$ on
weekly series to \ArimaMaxSlowdown{}$\times$ on \ArimaMaxGroup{}, since the
order search scales with series length and seasonal period while a forward pass
does not. On the \NumArimaGroups{} groups where it completed, it was beaten on
accuracy by the best pretrained model in \NumArimaBeaten{}. For a practitioner
the comparison that matters is therefore not among the pretrained models, whose
differences are mostly not resolvable, but between that family and the
automated classical search it would replace --- and there the case is strongest
exactly where the search is most expensive.

\section{Limitations}

The corpus-overlap finding is an association, not a demonstrated mechanism. It
rests on one corpus whose composition is publicly described and one domain that
matches it; a model could be better at Wikipedia traffic for reasons unrelated
to having read Wikipedia traffic, and our within-family comparison narrows that
possibility without closing it. The effect is also modest --- a rank-biserial
correlation of \CorpusEffect{} --- and we resist reading a large practical
consequence into it. Settling the question needs a benchmark stratified by
corpus membership across several models, which in turn needs disclosure that
most published corpora do not provide.

Beyond that: the hold-out is one forecast origin per group, where a
rolling-origin evaluation would estimate variance across time as well as across
series. Two domains contribute fewer than fifty series, so their rank intervals
are wide, and the electricity result should be read as ``no evidence of an
advantage'' rather than a demonstrated deficit. Pretraining cutoffs are taken
from model documentation and cannot be independently verified; where a cutoff
is undisclosed we substitute the release date. And the study covers five
domains, more than the standard two but far from the diversity of real
forecasting practice --- with the caveat, central to our argument, that adding
whichever domains are easiest to obtain is likely to add exactly the
well-published ones that pretraining corpora already contain.

\section{Conclusion}

Evaluated on data that postdates them, time-series foundation models are strong
on some domains, unnecessary on others, and useless where nothing works. We set
out to explain that spread with a property of the series and could not: neither
periodicity nor entropy orders it. What orders it is which domains the models
were raised on. The largest advantage in our study falls on the domain that
constitutes the bulk of one model's pretraining corpus, at the same
granularities, and it is that model's family which holds the advantage there
and nowhere else.

The practical consequence for benchmark design is that moving the test window
forward is not enough. It is necessary --- without it, nothing can be
concluded at all --- but it controls only for having seen these observations,
not for having spent a hundred billion time points learning what this kind of
observation looks like. A benchmark that wants to measure generalisation has to
hold out domains, not just dates, and it can only do so if pretraining corpora
are disclosed well enough to say which domains those are. The practitioner's
version of the same point is shorter: before asking which model is best, ask
whether your data looks like the data it was raised on.

We release the fetchers, the harness and the per-series losses so that both the
negative result and the corpus effect can be tested on domains we did not
consider.

\section*{Reproducibility}

Every dataset is rebuilt by a script in \texttt{data/sources/} requiring no
credentials. \texttt{scripts/run\_holdout.sh} reproduces the runs,
\texttt{scripts/significance.py} the tests, and \texttt{scripts/make\_tables.py}
every table and every number quoted above. Per-series losses are released
alongside the aggregates so that any alternative test can be applied without
re-running the models.

\bibliographystyle{plainnat}
\bibliography{refs}

\end{document}
